\documentclass[12pt,a4paper]{report}
\usepackage[utf8]{inputenc}
\usepackage[vietnamese]{babel}
\usepackage{amsmath, amsfonts, amssymb}
\usepackage{graphicx}
\usepackage{geometry}
\usepackage{booktabs}
\usepackage{enumitem}
\usepackage{setspace}
\usepackage{fancyhdr}
\usepackage{float}
\usepackage{listings}
\usepackage{xcolor}
\usepackage{hyperref}
\usepackage[most]{tcolorbox}
\usepackage{tikz}
\usetikzlibrary{shapes.geometric, arrows.meta, positioning, fit, backgrounds, shadow, calc}

% --- Cấu hình lề và khoảng cách ---
\geometry{margin=1in, top=1in, bottom=1in}
\onehalfspacing

% --- Cấu hình Code Snippets ---
\definecolor{codegreen}{rgb}{0,0.6,0}
\definecolor{codegray}{rgb}{0.5,0.5,0.5}
\definecolor{codepurple}{rgb}{0.58,0,0.82}
\definecolor{backcolour}{rgb}{0.95,0.95,0.92}

\lstset{
    backgroundcolor=\color{backcolour},   
    commentstyle=\color{codegreen},
    keywordstyle=\color{magenta},
    numberstyle=\tiny\color{codegray},
    stringstyle=\color{codepurple},
    basicstyle=\ttfamily\footnotesize,
    breaklines=true,                 
    captionpos=b,                    
    keepspaces=true,                 
    numbers=left,                    
    numbersep=5pt,                  
    tabsize=2,
    frame=lines
}

% --- Cấu hình Header/Footer ---
\pagestyle{fancy}
\fancyhf{}
\lhead{Báo cáo Tổng hợp Đồ án Tốt nghiệp}
\rhead{SEC-Wallet: Security-First Digital Wallet}
\cfoot{\thepage}

\begin{document}

% --- TRANG TIÊU ĐỀ ---
\begin{titlepage}
    \centering
    \vspace*{1cm}
    {\huge \textbf{KHOA CÔNG NGHỆ THÔNG TIN} \par}
    {\Large \textbf{CHUYÊN NGÀNH AN TOÀN THÔNG TIN} \par}
    \vspace{2cm}
    {\huge \textbf{BÁO CÁO TỔNG HỢP CUỐI KỲ} \par}
    \vspace{1cm}
    {\Huge \textbf{XÂY DỰNG HỆ THỐNG VÍ ĐIỆN TỬ BẢO MẬT \\ SEC-WALLET (V2.0)} \par}
    \vspace{2cm}
    \textbf{Sinh viên thực hiện:} \\
    Thái Hữu Thọ - 22127400 \\
    \vspace{1cm}
    \textbf{Giảng viên hướng dẫn:} \\
    Nguyễn Đình Thúc \\
    Ngô Đình Hy \\
    \vfill
    {\large Thành phố Hồ Chí Minh, \today \par}
\end{titlepage}

\tableofcontents
\newpage

\chapter{Phân tích Yêu cầu và Mục tiêu (GD1)}
\section{Bối cảnh dự án}
Trong kỷ nguyên số, các ứng dụng ví điện tử trở thành mục tiêu hàng đầu của tội phạm mạng. Đồ án **SEC-Wallet** tập trung vào việc hiện thực hóa các lớp bảo vệ mật mã học "quốc phòng" cho giao dịch tài chính cá nhân.

\section{Mục tiêu An ninh (CIA + N)}
Hệ thống được thiết kế để giải quyết 4 trụ cột an ninh:
\begin{itemize}
    \item \textbf{Confidentiality (Tính bí mật)}: Ngăn chặn nghe lén thông tin giao dịch bằng mã hóa AES-256-GCM.
    \item \textbf{Integrity (Tính toàn vẹn)}: Đảm bảo dữ liệu không bị thay đổi trên đường truyền thông qua Authentication Tag.
    \item \textbf{Availability (Tính khả dụng)}: Hệ thống hoạt động ổn định và chống tấn công từ chối dịch vụ (DoS) bằng Rate Limiting.
    \item \textbf{Non-repudiation (Chống chối bỏ)}: Ký số Ed25519 ràng buộc trách nhiệm pháp lý cho mỗi giao dịch.
\end{itemize}

\chapter{Kiến trúc Hệ thống và Cơ sở Dữ liệu (GD2 \& GD4)}

\section{Mô hình Thực thể Liên kết (ERD) Cập nhật}
Cơ sở dữ liệu SQLite (\texttt{sec\_wallet.db}) đã được mở rộng để hỗ trợ các tính năng bảo mật nâng cao.

\begin{figure}[H]
    \centering
    \resizebox{\textwidth}{!}{
    \begin{tikzpicture}[
        node distance=1.5cm and 1cm,
        table/.style={rectangle, draw=black, thick, fill=blue!5, text width=5cm, align=left, font=\small},
        header/.style={fill=blue!20, font=\bfseries\small},
        arrow/.style={-{Stealth[length=3mm]}, thick}
    ]
        % Nodes
        \node[table] (users) {
            \textbf{USERS} \\
            - id (PK) \\
            - username (U) \\
            - password\_hash \\
            - payment\_pin\_hash \\
            - \textbf{totp\_secret} \\
            - \textbf{failed\_login\_count} \\
            - salt
        };

        \node[table, right=of users] (wallets) {
            \textbf{WALLETS} \\
            - id (PK) \\
            - user\_id (FK) \\
            - encrypted\_balance \\
            - status
        };

        \node[table, below=of users, yshift=-1cm] (transactions) {
            \textbf{TRANSACTIONS} \\
            - id (PK) \\
            - sender\_id (FK-U) \\
            - receiver\_id (FK-U) \\
            - encrypted\_payload \\
            - auth\_tag \\
            - nonce \\
            - signature \\
            - \textbf{ephemeral\_pub} \\
            - \textbf{aad}
        };

        \node[table, right=of transactions] (logs) {
            \textbf{SECURITY\_LOGS} \\
            - id (PK) \\
            - event\_type \\
            - severity \\
            - \textbf{ip\_address} \\
            - actor\_user\_id (FK)
        };

        % Connections
        \draw[arrow] (users) -- (wallets);
        \draw[arrow] (users) -- (transactions);
        \draw[arrow] (users) -- (logs);
    \end{tikzpicture}
    }
    \caption{Mô hình ERD Cuối cùng của SEC-Wallet}
\end{figure}

\section{Sơ đồ Luồng Dữ liệu (DFD) Nâng cao}
Luồng dữ liệu tích hợp thêm bước **Xác thực 2 lớp (2FA)** và **Giám sát an ninh thời gian thực**.

\chapter{Thực thi Lõi Bảo mật (GD3 Implementation)}

\section{1. Mã hóa xác thực với AAD (AES-GCM-AAD)}
Hệ thống sử dụng trường \texttt{aad} để ràng buộc các metadata (như Transaction ID) vào gói tin mã hóa, giúp ngăn chặn việc hoán đổi gói tin giữa các phiên làm việc khác nhau.

\section{2. Perfect Forward Secrecy (PFS)}
Với trường \texttt{ephemeral\_pub} trong bảng giao dịch, hệ thống đảm bảo rằng nếu khóa static của người dùng bị lộ trong tương lai, các giao dịch cũ vẫn không thể bị giải băm vì mỗi lần thỏa thuận khóa đều sử dụng một cặp khóa tạm thời (Ephemeral Key).

\chapter{Tính năng Nâng cao (V2.0 Enhancements)}

\section{1. Xác thực 2 Lớp (TOTP)}
Hệ thống sử dụng tiêu chuẩn RFC 6238 để sinh mã OTP.
\begin{itemize}
    \item \textbf{Setup}: Người dùng quét mã QR để lưu \texttt{totp_secret}.
    \item \textbf{Verification}: Mọi giao dịch chuyển tiền đều yêu cầu OTP 6 chữ số từ Google Authenticator.
\end{itemize}

\section{2. Hệ thống Yêu cầu Thanh toán (Payment Requests)}
Tính năng cho phép khởi tạo yêu cầu nợ giữa các ví. 
\begin{itemize}
    \item \textbf{Request}: Tạo bản ghi Transaction với trạng thái \textit{Pending}.
    \item \textbf{Fulfill}: Người trả tiền nhập **Payment PIN** (Argon2) để ký số và thực thi chuyển tiền.
\end{itemize}

\chapter{Kiểm định và Bảo trì (GD5)}

\section{Kết quả Kiểm thử Đối kháng (security\_test.py)}
Script tự động đã xác thực 8 kịch bản tấn công:
\begin{table}[H]
    \centering
    \begin{tabular}{|l|l|l|}
    \hline
    \textbf{ID} & \textbf{Kịch bản tấn công} & \textbf{Kết quả} \\ \hline
    TEST 1 & Đăng nhập & JWT & PASS \\ \hline
    TEST 2 & Brute Force (Account Lockout) & PASS \\ \hline
    TEST 3 & Giao dịch hợp lệ (PFS) & PASS \\ \hline
    TEST 4 & Replay Attack (Nonce Store) & PASS \\ \hline
    TEST 5 & Tamper Attack (AES-GCM Auth) & PASS \\ \hline
    TEST 6 & Forgery Attack (EdDSA Sign) & PASS \\ \hline
    TEST 7 & JWT Forgery & PASS \\ \hline
    TEST 8 & Business Rule (Admin Protect) & PASS \\ \hline
    \end{tabular}
    \caption{Bảng Tổng hợp Kết quả Kiểm thử An ninh}
\end{table}

\section{Hoạt động Bảo trì Hệ thống}
- **Adaptive**: Di chuyển toàn bộ logic từ môi trường Local sang FastAPI Web.
- **Preventive**: Nâng cấp độ phức tạp của PBKDF2 lên mức an toàn tuyệt đối (600,000 vòng).

\chapter*{Kết luận và Hướng phát triển}
\addcontentsline{toc}{chapter}{Kết luận và Hướng phát triển}
Đồ án **SEC-Wallet** không chỉ dừng lại ở mức mô phỏng mà đã tiến sát tới các tiêu chuẩn của hệ thống ví điện tử thực tế về mặt an toàn mật mã. Trong tương lai, hệ thống có thể tích hợp thêm kiến trúc Microservices và cơ chế Multi-sig để tăng cường hơn nữa khả năng bảo vệ tài sản người dùng.

\end{document}
